\section{Wstęp}
Miara Lebesgue'a, czyli miary podzbiorów przestrzeni rzeczywistej $\R$ (i ogólniej, $\Rn$) z pracy Henriego Lebesgue’a, która dotyczyła rozszerzenia pojęcie całki na znacznie szerszą klasę funkcji niż całka Riemanna (tzw. całka Lebesgue’a). Teoria miary niemal od początku stawiała pytanie, czy istnieje miara, czyli funkcja przypisująca mierzalnym podzbiorom ustalonego zbioru nieujemne wartości liczbowe lub $+\infty$ w taki sposób, że zbiór pusty ma miarę zero, a miara sumy przeliczalnej rodziny parami rozłącznych zbiorów jest równa sumie ich miar. W teorii zbiorów aksjomatem wyboru taka miara nie istnieje. Miara Lebesgue'a spełnia drugi i trzeci z tych warnunków, a zamiast pierwszego podaje kryterium, które dzieli wszystkie możliwe zbiory zbioru liczb rzeczywistych na mierzalne oraz na niemierzalne. Definicja miary Lebesgue'a $\lambda$ opiera się na definicji miary zewnętrznej Lebesgue’a $\lambda^{*}$, która jest określona na zbiorze wszystkich podzbiorów przestrzeni rzeczywistej $\Rn$. Miara zewnętrzna Lebesgue’a jest niezmiennicza na przesunięcia oraz przedziałom jednostkowym przypisuje wartość 1, co znaczy, że jest odpowiednia do pomiaru wielkości geometrycznych, jak długość, pole powierzchni, objętość, itd. Miara ta nie jest jednak przeliczalnie addytywna, co oznacza, że miara sumy zbiorów rozłącznych może być mniejsza niż suma miar tych zbiorów. Właściwa miara Lebesgue'a nie ma już tej wady (tzn. jest przeliczalnie addytywna): jest równa mierze zewnętrznej Lebesgue'a, ale jej dziedzina jest ograniczona do zbiorów spełniających warunek Caratheodory’ego; zbiory te tworzą $\sigma$-ciało, które jest dziedziną miary Lebesgue’a. W pracy tej odtwarza się krok po kroku proces wiodący do konstrukcj tej miary.
